% !TEX program = lualatex
% Auto-generated from YAML (v3) — 08: データサイエンス・AI（1）：概要

\documentclass[aspectratio=43,professionalfonts,handout,10pt]{beamer}
\usepackage{beamer_template}

\newcommand{\LectureNum}{08}
\newcommand{\LectureTitle}{データサイエンス・AI（1）：概要}
\newcommand{\CourseName}{情報リテラシー}
\newcommand{\TextPages}{-}
\newcommand{\NextTopic}{データサイエンス・AI（2）：活用事例と身近な応用}
\newcommand{\NextPages}{pp.122-125}

\begin{document}
% --- Slide 1: title ---

\begin{frame}{\CourseName\quad 第\LectureNum 回「\LectureTitle 」}
  {\small \TermName\quad \TeacherName\quad ／\quad 教科書 \TextPages}

  \vfill

  \begin{block}{今日の流れ}
    \begin{enumerate}
      \item データサイエンスとは
      \item AI（人工知能）とは
      \item 身近なAI活用例
      \item Google Colabの使い方
      \item AI倫理の重要性
    \end{enumerate}
  \end{block}
  \vfill

  \begin{block}{今日の目標}
    \begin{itemize}
      \item データサイエンスやAI技術の概要を説明できる。
      \item これらを活用する際に必要なモラルや倫理を理解し，データを守るために必要な事項を説明できる。
    \end{itemize}
  \end{block}
\end{frame}

% --- Slide 2: twocol ---

\begin{frame}{データサイエンスとAIの違い}
  \begin{columns}[T]
    \begin{column}{0.48\textwidth}
      \begin{block}{データサイエンス}
        \small
        \begin{itemize}
          \item データから価値ある情報を引き出す技術
          \item 統計学・プログラミング・専門知識の融合
          \item 例：売上データから人気商品を見つける
          \item 例：センサーデータから設備の故障を予測する
          \item なぜ重要：データに基づいた意思決定・問題の早期発見・効率的な業務運営
        \end{itemize}
      \end{block}
    \end{column}
    \begin{column}{0.48\textwidth}
      \begin{exampleblock}{AI（人工知能）}
        \small
        \begin{itemize}
          \item 人間の知的活動をコンピュータで実現する技術
          \item 学習・推論・認識などの機能を持つ
          \item 例：スマホの顔認証、Siri・Alexa
          \item 例：Google翻訳、YouTubeのおすすめ動画
          \item データサイエンスの手法がAIの実現を支えている
        \end{itemize}
      \end{exampleblock}
    \end{column}
  \end{columns}
\end{frame}

% --- Slide 3: points ---

\begin{frame}{身近なAI活用例}
  \begin{block}{ポイント}
    \begin{itemize}
      \item スマホの顔認証 / YouTubeのおすすめ動画
      \item 自動翻訳（Google翻訳）/ スパムメールフィルター
      \item ゲームのAI / 音声入力（Siri・Alexa）
      \item カメラの美顔機能 / 検索エンジン
    \end{itemize}
  \end{block}

  \vfill

  \begin{exampleblock}{補足}
    \begin{itemize}
      \item 👉 気づかないところでたくさん使われている
    \end{itemize}
  \end{exampleblock}
\end{frame}

% --- Slide 4: terms ---

\begin{frame}{Google Colabとは}
  \begin{description}[labelwidth=6em]
    \item[\kw{Google Colaboratory（Colab）}]
      ブラウザ上でPythonを実行できる環境。インストール不要、Googleアカウントがあれば無料で使える
    \item[\kw{なぜColabを使うのか}]
      環境構築が不要・クラウドで動くのでPCのスペックに依存しない・データサイエンスでよく使われている
    \item[\kw{セルの実行方法}]
      Shift + Enter または左の▶ボタン
    \item[\kw{重要}]
      元のファイルは見るだけ（編集できない）。「ドライブにコピーを保存」をクリックして自分専用のファイルで作業する
  \end{description}
\end{frame}

% --- Slide 5: notice ---

\begin{frame}{AI倫理の重要性}
  \begin{noticebox}[注意]
    \begin{itemize}
      \item AIは道具：使い方次第で善にも悪にもなる。開発者・利用者の責任が重要
      \item なぜ倫理が必要か：社会への大きな影響／間違った使い方で人が傷つく
      \item 技術者として：技術だけでなく、その影響も考える倫理的判断力が必要
      \item 👉 工学を学ぶ者として、AIの社会的責任を意識しよう
    \end{itemize}
  \end{noticebox}
\end{frame}

% --- Slide 6: practice ---

\begin{frame}{Colab実習：Pythonで自己紹介グラフを作ろう}
  {\small 実際にコンピュータで操作してください。}

  \vfill

  \begin{block}{手順}
    \begin{enumerate}
      \item TeamsのチャットリンクからColabを開く
      \item 「ドライブにコピーを保存」をクリック
      \item 自己紹介を記入（名前・学科・趣味）
      \item グラフを作成：選択肢A「1週間の○○時間」／選択肢B「自由テーマ」
      \item データを変更して実行し、スクリーンショットを撮る
      \item Teamsの課題に提出
    \end{enumerate}
  \end{block}
\end{frame}

% --- チェックリスト ---

\begin{frame}{まとめ・チェックリスト}
  \begin{block}{自己チェック（できたら $\checkmark$ をつけよう）}
    \begin{itemize}
      \item[$\square$] データサイエンスとAIの概念と社会における役割を理解した
      \item[$\square$] 身近なAI活用例（画像認識・音声認識・推薦システム）を学んだ
      \item[$\square$] Google Colabの基本操作とAI倫理の重要性を把握した
    \end{itemize}
  \end{block}

  \vfill

  {\small 質問があれば授業後またはTeamsで受け付けます。}
\end{frame}

\end{document}
